% =============================================================================
%  text/HISTORICAL BACKGROUND/case-background.tex
%  Case background: demonstrates narrative paragraphs and a \citet citation.
%  案例背景示例。
% =============================================================================

% 请将此节替换为你自己的案例背景叙述。
Newly established regimes typically coalesce under leaders backed by cohorts
of capable elites. A robust literature establishes that elite capacity for
collective action is central to state-building \citep{author2018book,author2020example}.
Extensive evidence from developing regions, for instance, demonstrates that
elite fragmentation impairs rulers' ability to project authority and retards
the development of fiscal systems and public administration.

% 中文段落示例（可选）
% 新兴政体通常在由能干精英群体支持的领导者领导下凝聚起来。大量研究证实，
% 精英阶层的集体行动能力对于国家建设至关重要；精英在地理或派系线上的
% 碎片化会削弱统治者推行权威的能力，并阻碍财政体系与行政机构的发展。

Yet rulers' rational calculations do not invariably favor elite cohesion.
As \citet{author2019chapter} argues, the strategic imperatives of political
survival may diverge sharply from the requirements of national prosperity.
History bears this out: many prominent rulers fell not to external adversaries
or mass uprisings, but to challenges orchestrated by their closest associates.

% 历史背景部分的其余内容（事件梳理、关键时间线等）请在此展开。
Herein lies the fundamental tension this paper studies: an elite network
powerful enough to sustain the state simultaneously poses a latent threat to
the ruler's own position. The case examined below offers sharp temporal
variation in the ruler's power security, enabling us to observe how the same
ruler adopted diametrically opposed strategies under different conditions.

Chinese-language scholarship makes closely related arguments about the
relationship between institutional change and state capacity
\citep{zhang2018example}; engaging this literature is encouraged when
relevant to your case.
